\documentclass[12pt,a4paper]{article}

% --- Настройка языка и шрифтов для LuaLaTeX ---
\usepackage{fontspec}
\defaultfontfeatures{Ligatures=TeX,Scale=MatchLowercase}

\setmainfont{Alice-Regular}[
    Path=font/,
    Extension=.ttf,
    BoldFont = Alice-Regular,
    ItalicFont = Alice-Regular,
    BoldItalicFont = Alice-Regular,
    BoldFeatures = {FakeBold=1.6},
    ItalicFeatures = {FakeSlant=0.2},
    BoldItalicFeatures = {FakeBold=1.6,FakeSlant=0.2}
]

\newfontfamily\cyrillicfont{Alice-Regular}[
    Path=font/,
    Extension=.ttf,
    BoldFont = Alice-Regular,
    ItalicFont = Alice-Regular,
    BoldItalicFont = Alice-Regular,
    BoldFeatures = {FakeBold=1.6},
    ItalicFeatures = {FakeSlant=0.2},
    BoldItalicFeatures = {FakeBold=1.6,FakeSlant=0.2}
]

\usepackage{polyglossia}
\setmainlanguage{russian}
\setsansfont{TeX Gyre Heros}
\setmonofont{TeX Gyre Cursor}

\usepackage{graphicx}        					% Графика
\usepackage{amsmath, amsfonts, amssymb, amsthm}	% Математика
\usepackage{mathtools}
\usepackage{chemfig}
\usepackage[version=4,arrows=pgf-filled,
            textfontname=sffamily,
            mathfontname=mathsf]{mhchem}
\usepackage[protrusion=false]{microtype}       % Микротипографика
\usepackage{braket} 						% Braket нотация для квантовой механики
\usepackage{unicode-math} 					% Современный пакет для математики в LuaLaTeX
\usepackage{longtable}       				% Таблицы с разрывом страниц
\usepackage{tikz}           				% Рисование
\usepackage[europeanresistors]{circuitikz}	% Рисование электрических схем в европейском стиле
\usepackage{listings}						% Для отображения кода
\usepackage{tabularx}       				% Для таблиц во всю ширину
\usepackage{array}		  					% Для расширенных возможностей таблиц
\usepackage{float}          				% Для точного позиционирования фигур
\usepackage{pgfplots}                          % Для построения графиков
\usepackage{caption}                           % Для настройки подписей к рисункам и таблицам
\usepackage{xfp}                               % Для вычислений в LaTeX
\usepackage{luacode}                           % Для выполнения Lua-кода в документе    
\pgfplotsset{compat=1.18}
\usepackage[paper=letterpaper,
            top=20mm,
            bottom=20mm,
            left=10mm,
            right=10mm,
            footskip=0.4in 
]{geometry}

\usetikzlibrary{arrows.meta}

% --- Колонтитулы ---
\usepackage{fancyhdr}
\pagestyle{fancy}
\fancyhf{}

% --- Теоремы и окружения ---
\theoremstyle{plain}
\newtheorem{theorem}{Theorem}[section]
\newtheorem{lemma}[theorem]{Lemma}
\newtheorem{proposition}[theorem]{Proposition}
\newtheorem{corollary}[theorem]{Corollary}

\theoremstyle{definition}
\newtheorem{definition}[theorem]{Definition}

\theoremstyle{remark}
\newtheorem{remark}[theorem]{Remark}

\numberwithin{equation}{section}


\renewcommand{\headrulewidth}{0.4pt}
\renewcommand{\footrulewidth}{0.4pt}
\renewcommand\lstlistingname{Исходный код}

\lstdefinestyle{main}{
	backgroundcolor=\color{gray!8},   
	commentstyle=\color{green!50!black},
	keywordstyle=\color{blue!80!black},
	numberstyle=\small\color{darkgray},
	stringstyle=\rmfamily\color{orange!90!black},
	basicstyle=\ttfamily\small,
	breakatwhitespace=false,         
	breaklines=true,                   
	keepspaces=true,                 
	numbers=left,                    
	numbersep=8pt,                  
	showspaces=false,                
	showstringspaces=false,
	showtabs=false,                  
	tabsize=4
}	

\lstset{style=main}

\newcolumntype{C}{>{\centering\arraybackslash}X}
\renewcommand{\arraystretch}{1.25}

\begin{luacode*}
    data = require("data")
    include = require("include")
    handle = require("handle")
\end{luacode*}

\begin{document}

\section{Построить прямую ветвь ВАХ германиевого и кремниевого диодов в 
полулогарифмических координатах при температурах $T = 20^{\circ}\text{C}$ 
и $T = 75^{\circ}\text{C}$. Определить параметры диодов $m$ и $I_0$.}

\centering\directlua{
    local points_D1_20Celsius_x = data.table1[1]
    local points_D1_20Celsius_y = data.table1[2]
    local table1_D1_20Celsius = {
        "yellow!30!gray",
        "*",
        "2pt",
        [[Кремниeвый диод при $T = 20^{\circ}\mathrm{C}$]],
        points_D1_20Celsius_x,
        points_D1_20Celsius_y
    }
    local points_D1_75Celsius_x = data.table1[1]
    local points_D1_75Celsius_y = data.table1[3]
    local table1_D1_75Celsius = {
        "red!30!gray",
        "*",
        "2pt",
        [[Кремниeвый диод при $T = 75^{\circ}\mathrm{C}$]],
        points_D1_75Celsius_x,
        points_D1_75Celsius_y
    }
    local points_D2_20Celsius_x = data.table1[1]
    local points_D2_20Celsius_y = data.table1[4]
    local table1_D2_20Celsius = {
        "yellow!30!black",
        "*",
        "2pt",
        [[Германиевый диод при $T = 20^{\circ}\mathrm{C}$]],
        points_D2_20Celsius_x,
        points_D2_20Celsius_y
    }
    local points_D2_75Celsius_x = data.table1[1]
    local points_D2_75Celsius_y = data.table1[5]
    local table1_D2_75Celsius = {
        "red!30!black",
        "*",
        "2pt",
        [[Германиевый диод при $T = 75^{\circ}\mathrm{C}$]],
        points_D2_75Celsius_x,
        points_D2_75Celsius_y
    }

    include.pgf_plot_wrap("Прямая ветвь ВАХ германиевого и кремниевого диодов", 
                          "x", "y", 
                          0.9, 0.5, table1_D1_20Celsius, 
                                    table1_D1_75Celsius, 
                                    table1_D2_20Celsius, 
                                    table1_D2_75Celsius)
}

\section{Построить прямую ветвь ВАХ германиевого и кремниевого диодов в 
линейных координатах при $T = 20^{\circ}\text{C}$ и $T = 75^{\circ}\text{C}$. 
Определить объемное сопротивление базы $r_\text{б}$ и ток вырождения $I_\text{в}$ для каждой ВАХ. }

\section{Вычислить дифференциальное сопротивление диодов $r_d$ при токах $I = 1$ мА и $I = 10$ мА.}

\section{Определить температурный коэффициент напряжения диодов при токах $10$ мкА, $1$ мА, $20$ мА.}

\section{Построить обратную ветвь ВАХ диодов в линейных координатах при $T = 20^{\circ}\text{C}$ и 
$T = 75^{\circ}\text{C}$. Определить характер изменения обратного тока диодов в 
зависимости от приложенного напряжения.}

\end{document}
